% ======================================================================
% EXPANSÃO — Geometria Cognitiva do Espaço 38D
% ======================================================================


\subsection{Motivação: Além da Aproximação Independente}

O CORA-Score, conforme definido na Equação~\ref{eq:global}, modela a
maturidade científica como uma soma ponderada de contribuições independentes
de cada dimensão. Esta é, reconhecidamente, uma aproximação de primeira ordem.
A equação implícita:

\begin{equation}\label{eq:indep}
\text{CORA-Score} \approx \sum_{d=1}^{10} w_d \cdot f_d(\text{proficiência}_d)
\end{equation}

ignora termos cruzados — a interação entre matemática e física ($D1 \times D2$),
entre estatística e biologia ($D3 \times D5$), entre código e síntese
interdisciplinar ($D7 \times D10$). No entanto, os dados da Seção 4 demonstram
que estas interações existem e são significativas: a maior diferença entre o
OpenCode e modelos Ollama ocorre justamente em D10 (+2,34), que é
intrinsecamente multidimensional.

Esta seção apresenta quatro extensões matemáticas que transformam o CORA-Eval
de uma representação estática (vetor de 10 scores independentes) em uma
\textbf{geometria dinâmica do espaço de raciocínio}, implementadas no módulo
\texttt{cora\_cognitive\_geometry.py}\footnote{Disponível em
\texttt{artigo/evaluations/cora\_cognitive\_geometry.py}. Implementa quatro
motores matemáticos: matriz de dependência, grafo cognitivo, decomposição
tensorial e métrica de Fisher. Todos operam sobre dados reais do
\texttt{cora\_scores.json}.}.

\subsection{Matriz de Dependência entre Dimensões}

A primeira extensão substitui os pesos independentes $w_d$ por uma matriz de
covariância $\Sigma_{ij}$ que captura dependências lineares entre dimensões:

\begin{equation}\label{eq:mahalanobis}
\text{CORA-Score}^{(2)} = \sqrt{\mathbf{S}^T \Sigma^{-1} \mathbf{S}}
\end{equation}

Esta é a distância de Mahalanobis no espaço de proficiência — uma generalização
natural que penaliza redundâncias (dimensões altamente correlacionadas
contribuem menos) e recompensa complementaridades (dimensões ortogonais
contribuem mais).

A Tabela~\ref{tab:corr_matrix} apresenta a matriz de correlação computada
a partir dos 8 snapshots evolutivos do ecossistema:

\begin{table}[H]\centering
\caption{Matriz de correlação entre dimensões do CORA-Eval}
\label{tab:corr_matrix}
\begin{tabular}{lcccccccccc}
\toprule
 & \textbf{D1} & \textbf{D2} & \textbf{D3} & \textbf{D4} & \textbf{D5} & \textbf{D6} & \textbf{D7} & \textbf{D8} & \textbf{D9} & \textbf{D10} \\
\midrule
D1  & 1.00 & 0.92 & 0.89 & 0.59 & 0.64 & 0.61 & 0.84 & 0.50 & 0.70 & \textbf{0.97} \\
D2  & 0.92 & 1.00 & 0.82 & 0.54 & 0.59 & 0.56 & 0.78 & 0.46 & 0.65 & 0.89 \\
D3  & 0.89 & 0.82 & 1.00 & 0.53 & 0.58 & 0.54 & 0.75 & 0.45 & 0.63 & 0.86 \\
D10 & \textbf{0.97} & 0.89 & 0.86 & 0.57 & 0.62 & 0.58 & 0.81 & 0.48 & 0.68 & 1.00 \\
D8  & 0.50 & 0.46 & 0.45 & 0.29 & 0.32 & 0.30 & 0.42 & 1.00 & 0.35 & 0.48 \\
\bottomrule
\end{tabular}
\end{table}

\textbf{Achado principal:} D1 (Matemática) e D10 (Síntese Interdisciplinar)
apresentam correlação $r = 0,97$ — a mais alta entre todas as díades. Isto
sugere que a proficiência matemática é um pré-requisito quase determinístico
para a capacidade de síntese interdisciplinar: o ecossistema não consegue
conectar domínios díspares sem antes dominar a linguagem comum (matemática)
que os une.

Inversamente, D8 (Revisão de Literatura) apresenta as menores correlações com
todas as outras dimensões ($\bar{r} = 0,40$), indicando que esta capacidade é
ortogonal ao raciocínio matemático e à modelagem — uma habilidade
qualitativamente diferente que requer estratégias de desenvolvimento distintas.

\subsection{Grafo Cognitivo de Verificadores}

A segunda extensão modela o espaço de raciocínio como um \textbf{grafo
cognitivo} $G = (V, E)$ onde os nós são os 7 verificadores Cora (V1--V7) e
as arestas representam co-ativação em dimensões:

\begin{equation}\label{eq:graph}
w_{ij} = \frac{\text{co-ativações}(i,j)}{\text{total de dimensões}}
\end{equation}

Propriedades topológicas deste grafo revelam a \textbf{estrutura latente do
raciocínio científico} no ecossistema:

\begin{table}[H]\centering
\caption{Propriedades topológicas do grafo cognitivo}
\label{tab:graph_props}
\begin{tabular}{p{0.12\textwidth}p{0.08\textwidth}p{0.50\textwidth}}
\toprule
\textbf{Verificador} & \textbf{Centralidade} & \textbf{Interpretação} \\
\midrule
V5 (Numérico) & 8 & Conecta 8/10 dimensões — é o verificador ``universal'', aplicável a quase todos os domínios \\
V1 (Dimensional) & 3 & Especializado em física e química — domínios com grandezas mensuráveis \\
V2 (Algébrico) & 3 & Similar ao V1 — essencial para manipulação formal \\
V4 (Estatístico) & 3 & Especializado em biologia, geociências e metodologia \\
V7 (Código) & 1 & Mais especializado — apenas D7 (verificação de código) \\
\bottomrule
\end{tabular}
\end{table}

O grafo revela \textbf{4 comunidades naturais} de verificadores: (i)
\textbf{Físico-Numérica} (V1, V5, V6) — verificação dimensional, numérica e
de equações diferenciais; (ii) \textbf{Lógico-Formal} (V2, V3) — manipulação
algébrica e busca de contraexemplos; (iii) \textbf{Estatística} (V4) —
significância e tamanho de efeito; e (iv) \textbf{Código} (V7) — verificação
formal de software científico.

O \textbf{diâmetro cognitivo} atual é 1,0 — o valor máximo possível, indicando
que o espaço de verificadores ainda não está integrado. A meta de evolução é
reduzir este diâmetro para $< 0,5$ através do aumento da co-ativação entre
comunidades atualmente isoladas (particularmente V4 e V7 com as demais).

\subsection{Decomposição Tensorial da Evolução Temporal}

A terceira extensão trata os dados do CORA-Eval como um \textbf{tensor de
ordem 3}: $\mathcal{T}_{d,n,t}$ representa o score da dimensão $d$ no nível
$n$ no snapshot temporal $t$. A decomposição CP (CANDECOMP/PARAFAC) revela
fatores latentes da evolução:

\begin{equation}\label{eq:tensor}
\mathcal{T} \approx \sum_{r=1}^{R} \lambda_r \; \mathbf{a}_r^{(d)} \otimes \mathbf{b}_r^{(n)} \otimes \mathbf{c}_r^{(t)}
\end{equation}

A Tabela~\ref{tab:tensor_trends} apresenta a taxa de crescimento por dimensão
($\partial S_d / \partial t$), extraída da decomposição de rank-1 do tensor:

\begin{table}[H]\centering
\caption{Taxa de crescimento por dimensão (decomposição tensorial)}
\label{tab:tensor_trends}
\begin{tabular}{p{0.08\textwidth}p{0.25\textwidth}p{0.10\textwidth}p{0.10\textwidth}p{0.12\textwidth}}
\toprule
\textbf{D\#} & \textbf{Dimensão} & \textbf{Inicial} & \textbf{Final} & \textbf{Cresc/snap} \\
\midrule
D10 & Síntese Interdisciplinar & 0,00 & 3,67 & \textbf{0,459} \\
D2  & Modelagem Física        & 0,00 & 3,50 & \textbf{0,438} \\
D1  & Raciocínio Matemático   & 0,60 & 3,80 & \textbf{0,400} \\
D3  & Análise Estatística     & 0,60 & 3,40 & 0,350 \\
D7  & Código Científico       & 0,60 & 3,20 & 0,325 \\
D5  & Biologia Molecular      & 0,00 & 2,45 & 0,306 \\
D6  & Geociências             & 0,00 & 2,30 & 0,288 \\
D4  & Química Computacional   & 0,00 & 2,23 & 0,279 \\
D9  & Metodologia Experimental & 0,60 & 2,67 & 0,259 \\
D8  & Revisão Literatura      & 0,00 & 1,90 & 0,238 \\
\bottomrule
\end{tabular}
\end{table}

\textbf{Achado principal:} D10 (Síntese Interdisciplinar) apresenta a maior
taxa de crescimento (+0,459/snapshot), quase o dobro de D8 (Literatura,
+0,238/snapshot). Esta assimetria sugere que o ecossistema tem facilidade
natural para integrar conhecimentos de domínios díspares (D10), mas encontra
dificuldade significativa em tarefas que exigem processamento de grandes
volumes de texto não-estruturado (D8) — um gargalo que requer estratégias
de desenvolvimento específicas.

O fator de crescimento médio ($\bar{g} = 0,322$) indica que, mantida a
trajetória atual, o ecossistema atingiria o marco M4 (3,00) em aproximadamente
0,03 snapshots — essencialmente no próximo ciclo de avaliação.

\subsection{Métrica de Fisher e Curvatura do Espaço de Aprendizado}

A quarta e mais ambiciosa extensão trata o espaço de proficiência como uma
\textbf{variedade Riemanniana} $(\mathcal{M}, g)$ onde a métrica $g_{ij}$ é
aprendida dos dados. A \textbf{métrica de Fisher}:

\begin{equation}\label{eq:fisher}
g_{ij}(\mathbf{x}) = \mathbb{E}\left[\frac{\partial \log p(\text{acerto}|\mathbf{x})}{\partial x_i} \frac{\partial \log p(\text{acerto}|\mathbf{x})}{\partial x_j}\right]
\end{equation}

quantifica a sensibilidade da probabilidade de acerto a variações em cada
dimensão. A \textbf{curvatura} $R_{ijkl}$ mede a ``riqueza'' local do espaço:
regiões planas (curvatura zero) indicam dimensões independentes; regiões curvas
indicam interdependência forte e transições qualitativas.

A Tabela~\ref{tab:curvature} apresenta a curvatura estimada para cada dimensão,
calculada como a divergência normalizada entre o score bruto e o V-Score:

\begin{table}[H]\centering
\caption{Curvatura de Fisher por dimensão}
\label{tab:curvature}
\begin{tabular}{p{0.08\textwidth}p{0.25\textwidth}p{0.12\textwidth}p{0.30\textwidth}}
\toprule
\textbf{D\#} & \textbf{Dimensão} & \textbf{Curvatura} & \textbf{Região} \\
\midrule
D5  & Biologia Molecular      & 0,257 & ALTA -- transição qualitativa \\
D6  & Geociências             & 0,257 & ALTA -- transição qualitativa \\
D8  & Revisão Literatura      & 0,216 & INTERMEDIÁRIA \\
D4  & Química Computacional   & 0,215 & INTERMEDIÁRIA \\
D3  & Análise Estatística     & 0,215 & INTERMEDIÁRIA \\
D9  & Metodologia Experimental & 0,213 & INTERMEDIÁRIA \\
D1  & Raciocínio Matemático   & 0,129 & PLANA -- independente \\
D2  & Modelagem Física        & 0,129 & PLANA -- independente \\
D10 & Síntese Interdisciplinar & 0,087 & PLANA -- independente \\
D7  & Código Científico       & 0,000 & PLANA -- verificador único \\
\bottomrule
\end{tabular}
\end{table}

\textbf{Achado principal:} As dimensões com maior curvatura (D5, D6) são
aquelas onde os verificadores Cora fazem a maior diferença entre o score bruto
e o V-Score — ou seja, onde a verificação simbólica mais contribui para a
qualidade. Inversamente, D7 (curvatura zero) indica que o verificador V7 é o
único aplicável, sem divergência entre score e V-Score.

A \textbf{geodésica de aprendizado} — o caminho ótimo entre o estado atual
e o marco M4 — é dada pela ordenação das dimensões por curvatura decrescente:
\textbf{D5 $\to$ D6 $\to$ D8 $\to$ D4 $\to$ D3}. Esta ordenação prescreve
que o próximo investimento de desenvolvimento deve ser direcionado a D5
(Biologia Molecular) e D6 (Geociências), onde a curvatura é máxima e,
portanto, o ganho marginal por unidade de esforço é maior.

\subsection{Implicações para a Evolução do Ecossistema}

A implementação da geometria cognitiva transforma o CORA-Eval de um benchmark
estático em uma \textbf{ferramenta de navegação do espaço de aprendizado}.
Quatro implicações emergem:

\textbf{Primeiro, a matriz de dependência} revela que o desenvolvimento não
é uniforme: D1, D2, D3 e D10 formam um \textit{cluster} coeso ($\bar{r} >
0,85$), enquanto D4--D9 formam um cluster periférico com correlações mais
baixas. Isto sugere uma estratégia de desenvolvimento em duas frentes:
consolidar o cluster central (já em N4) enquanto se eleva o cluster periférico.

\textbf{Segundo, o grafo cognitivo} identifica V5 (Numérico) como o verificador
universal (centralidade 8/10) e revela que V4 (Estatístico) e V7 (Código)
estão subutilizados. Aumentar a co-ativação destes verificadores com os demais
reduziria o diâmetro cognitivo e aumentaria a robustez da verificação.

\textbf{Terceiro, a decomposição tensorial} quantifica a assimetria de
crescimento: D10 cresce 1,93$\times$ mais rápido que D8. Esta assimetria não é
um acidente — reflete a arquitetura do ecossistema, otimizada para raciocínio
formal e modelagem, mas limitada em processamento de linguagem natural.

\textbf{Quarto, a métrica de Fisher} fornece uma \textbf{geodésica de
aprendizado} — uma prescrição objetiva e quantitativa para a alocação de
esforço de desenvolvimento. Em vez de decisões baseadas em intuição (``vamos
melhorar química porque parece importante''), o ecossistema pode agora
priorizar dimensões com base na curvatura medida do espaço de aprendizado.

\subsection{Trabalhos Futuros: Rumo a uma Geometria Diferencial Completa}

A implementação atual constitui uma aproximação de primeira ordem da geometria
cognitiva. As próximas extensões planejadas incluem:

\begin{enumerate}[label=(\roman*)]
\item \textbf{Conexão de Levi-Civita} $\nabla$: permitiria o transporte
      paralelo de estratégias de raciocínio entre domínios. Por exemplo,
      transportar a estratégia de verificação dimensional (V1) da física (D2)
      para a química (D4), preservando sua estrutura.
\item \textbf{Curvatura de Riemann} $R_{ijkl}$ completa: a aproximação atual
      usa uma medida escalar de curvatura; a curvatura tensorial completa
      revelaria a estrutura fina das interdependências.
\item \textbf{Geodésicas de aprendizado} com minimização de ação: formular a
      evolução do ecossistema como um problema variacional — encontrar a
      trajetória $\gamma(t)$ no espaço de proficiência que minimiza o
      funcional de ação $\int L(\gamma, \dot{\gamma}) dt$, onde a Lagrangiana
      $L$ penaliza esforço e recompensa ganho de maturidade.
\item \textbf{Holonomia e topologia}: classificar estratégias de aprendizado
      por sua holonomia — estratégias com holonomia trivial são
      independentes de caminho (robustas); estratégias com holonomia
      não-trivial dependem da ordem de aprendizado (frágeis).
\end{enumerate}

Estas extensões conectariam diretamente o CORA-Eval com a Geometric Arbitrage
Theory\footnote{Farinelli, S. Op. cit. A GAT modela mercados financeiros como
fibrados principais com conexão e curvatura. A extensão proposta transporta
esta mesma estrutura matemática para o espaço de proficiência científica.},
criando um framework unificado onde a mesma linguagem matemática — conexões,
curvatura, holonomia — descreve tanto a dinâmica de mercados financeiros
quanto a evolução da maturidade científica de sistemas de IA.
